\documentclass[11pt,letterpaper]{article}

\usepackage[utf8]{inputenc}
\usepackage[T1]{fontenc}
\usepackage[spanish,es-nodecimaldot]{babel}
\usepackage{lmodern}
\usepackage[margin=2.3cm]{geometry}
\usepackage{amsmath,amssymb}
\usepackage{booktabs}
\usepackage{tabularx}
\usepackage{graphicx}
\usepackage{float}
\usepackage{xcolor}
\usepackage{hyperref}

\hypersetup{
    colorlinks=true,
    linkcolor=blue!50!black,
    urlcolor=blue!50!black
}

\setlength{\parindent}{0pt}
\setlength{\parskip}{0.45em}

\title{Modelación de la cadena de suministro de ayuda humanitaria\\
\large Ciudad UVG: terremoto de magnitud 6.8}
\author{Grupo 3}
\date{\today}

\begin{document}

\maketitle

\section{Justificación del paradigma}
% Extensión objetivo: 1 página.
% Comparar SD, DES y ABM y justificar el paradigma seleccionado.

\section{Descripción del modelo}
% Extensión objetivo: 2 páginas.
% Presentar el ODD simplificado: entidades, atributos, ecuaciones,
% scheduling, comunicación y distribución inicial de atributos.

\subsection{Propósito y alcance}

\subsection{Entidades y variables de estado}

\subsection{Ecuaciones y reglas de operación}

\subsection{Scheduling e inicialización}

\section{Resultados y análisis}
% Extensión objetivo: 3 páginas.
% Incluir trayectorias, medias, intervalos de confianza del 95 %,
% momentos críticos y los outputs requeridos por el Excel.

\subsection{Balance de suministros}

\subsection{Primer agotamiento crítico y margen de respuesta}

\subsection{Recomendación logística}

\section{Incorporación del intercambio}
% Extensión objetivo: 1 página.
% Documentar lo recibido de los Grupos 1 y 4, su incorporación y
% los cambios producidos en las conclusiones.

\section{Limitaciones y propuestas de mejora}
% Extensión objetivo: 1 página.
% Incluir al menos dos limitaciones estructurales y sus mejoras.